% !LOOM digest: Arden24
% !LOOM prefix: Arden24
% !LOOM extracted-from: arXiv:2504.01234v1
% !LOOM method: extract
% !LOOM created: 2026-09-14
% !LOOM numbering: emulated

\section*{Overview}
Let $Q$ be a finite quiver. A weighting of $Q$ assigns an integer to each arrow, and the weightings form a free abelian group of rank $|A(Q)|$. Among them sit the weightings that balance at every vertex. This note fixes the notation for that sublattice and proves the three facts about it that are used everywhere.

Throughout, all quivers are finite, loops and parallel arrows are permitted, and the underlying graph of a quiver is the undirected multigraph obtained by forgetting the orientation of each arrow.

\begin{theorem}[{\cite[Standing assumptions]{Arden24}}]\label{Arden24-setup}
\emph{Sentences of the paper that state an assumption, gathered by the extractor; check each one's scope.}

Throughout, all quivers are finite, loops and parallel arrows are permitted, and the underlying graph of a quiver is the undirected multigraph obtained by forgetting the orientation of each arrow.
\end{theorem}

\section{Introduction}\label{Arden24-sec-1}

\begin{definition}[{\cite[Definition 1.1, p.~1]{Arden24}}]\label{Arden24-def-1.1}
\label{Arden24-def:quiver}
A \emph{quiver} is a quadruple $Q = (V, A, s, t)$ in which $V$ and $A$ are finite sets and $s, t \colon A \to V$ are functions, called the source and target maps.
\end{definition}

\begin{definition}[{\cite[Definition 1.2, p.~1]{Arden24}}]\label{Arden24-def-1.2}
\label{Arden24-def:cycle-space}
The \emph{cycle space} of a quiver $Q$ is the kernel $\mathcal{Z}(Q)$ of the divergence map $\partial \colon \mathbb{Z}^{A(Q)} \to \mathbb{Z}^{V(Q)}$ whose value at a vertex $v$ is the total weight of the arrows with target $v$ less the total weight of the arrows with source $v$.
\end{definition}

\section{The image of the divergence map}\label{Arden24-sec-2}

\begin{proposition}[{\cite[Proposition 2.1, p.~1]{Arden24}}]\label{Arden24-prop-2.1}
\label{Arden24-prop:image}
Let $Q$ be a quiver whose underlying graph has $c$ connected components. The image of the divergence map is the lattice of integer vectors in $\mathbb{Z}^{V(Q)}$ whose coordinates sum to zero on each component, and its rank is $|V(Q)| - c$.
\end{proposition}

\begin{lemma}[{\cite[Lemma 2.2, p.~2]{Arden24}}]\label{Arden24-lem-2.2}
\label{Arden24-lem:saturated}
The cycle space of a quiver is a saturated sublattice of the lattice of weightings: if $nw$ is balanced for some nonzero integer $n$, then $w$ is balanced.
\end{lemma}

\section{Rank and duality}\label{Arden24-sec-3}

\begin{theorem}[{\cite[Theorem 3.1, p.~2]{Arden24}}]\label{Arden24-thm-3.1}
\uses{Arden24-prop-2.1}
\label{Arden24-thm:rank}
Let $Q$ be a quiver whose underlying graph has $c$ connected components. Then the cycle space of $Q$ is a free abelian group of rank $|A(Q)| - |V(Q)| + c$, and this rank depends only on the underlying graph of $Q$.
\end{theorem}

\begin{corollary}[{\cite[Corollary 3.2, p.~2]{Arden24}}]\label{Arden24-cor-3.2}
\uses{Arden24-thm-3.1, Arden24-lem-2.2}
\label{Arden24-cor:forest}
The cycle space of a quiver is trivial if and only if the underlying graph is a forest.
\end{corollary}

\section{Examples and remarks}\label{Arden24-sec-4}

\begin{remark}[{\cite[Remark 4.1, p.~2]{Arden24}}]\label{Arden24-rem-4.1}
\label{Arden24-rem:orientation}
Nothing above uses the orientation of the arrows except as bookkeeping, which is the content of the last sentence of the proof of Theorem~\ref{Arden24-thm-3.1}. Readers who prefer graphs to quivers may replace every quiver by its underlying graph with an arbitrary orientation, and no statement changes.
\end{remark}

\begin{remark}[{\cite[Remark 4.2, p.~2]{Arden24}}]\label{Arden24-rem-4.2}
\label{Arden24-rem:infinite}
For an infinite quiver of finite valence the divergence map is still defined on finitely supported weightings and its kernel is still saturated, but the rank formula of Theorem~\ref{Arden24-thm-3.1} has no meaning, since all three of its terms are infinite. The correct statement in that setting is a description of the cycle space as a direct limit, which we do not pursue.
\end{remark}
